\section{怎样辨析文言介词的一般用法}

古代汉语的介词，同现代汉语的介词一样，其语法作用都是同名词、代词或名词性词组组成介词结构，用来修饰动词或形容词，表示动作行为的处所、时间、对象、方式或形成的原因、目的等，有时引进动作行为的主动者或比较的对象。不同的是，古代汉语介词数量少，主要有“以、为、于、与、因”（另外还有“乎、自、由、从、在、至、见、用、被、当”等），因此，一个文言介词在不同语言环境中能相当现代汉语不少介词的讲法，这是一；由于古代汉语的文字少，所以一个文言介词往往身兼几类词的用法，这是二；还由于古今词义的发展变化，使得某些介词（如“因、见”等）的古今意义差别很大，这是三。正因为有这些不同，我们要掌握好文言介词的用法就有一定的困难。好在古今汉语的介词有很多共同点，我们如果再掌握了辨析文言介词的基本方法，就能准确地理解文言介词在文句中的意义和作用。我们现在就以几个常见的文言介词“以、为、于、因、与”为例，来谈谈辨析文言介词的基本方法。

根据句意和介词特点来确定介词的词性。

这是辨析文言介词的基本出发点。因为如果不能确定是不是介词，那么对它的意义和语法作用就无从谈起。

由于文言虚词中的纯粹介词很少，绝大部分都是一词身兼几类词性（例如“以、因、与”等，既作介词用，又可作连词用；“以、为、从、在、至、见”等介词，本身原是动词，仍可作动词用；“乎、与、为”等介词，有时又可作语气词或助词用），因此，首先要辨清词性，才能进一步辨析其词义或作用。
辨析一个词是不是介词的基本方法，一是根据句意，二是根据介词特点。辨别介宾结构时（或叫介词结构，介词词组），更应注意分清句中的主要动词（或用作动词的形容词），分清动宾与介宾的不同。例如：

[以]
①敌人远我，欲
以火器
以火器困我也。（《冯婉贞》）
②而吾
以捕蛇
以捕蛇独存。（《捕蛇者说》）
③余船
以次
以次俱进。（《赤壁之战》）
④余以乾隆三十九年十二月，自京师乘风雪，历齐河、长清……至于泰安。（《登泰山记》）
⑤夷以近，则游者众；险以远，则至者少。（《游褒禅山记》）
⑥属予作文以记之。（《岳阳楼记》）
⑦臣之妻私臣，臣之妾畏臣，臣之客欲有求于臣，皆以美于徐公。（《邹忌讽齐王纳谏》）
我们可以根据句意和介词特点辨明①②③④句的“以”都是介词，这些介词后面都有句中的主要动词。（⑤⑥句的“以”是连词，⑦句的“以”是动词。）
①的“以”与名词“火器”构成介宾结构，修饰谓语动词“困”。这类介词“以”字是表示动作行为以某物为工具或凭借，相当于现代的“用、拿、把、凭”。
②的“以”与动宾词组“捕蛇”构成介宾结构，修饰谓语动词“存”。这类介词“以”字是表示原因，相当于现代的“因、因为、由于”。
③的“以”与名词“次”构成介宾结构，修饰主动词“进”。这类介词“以”是表示“依照”、“按着”的意思。

④的“以”与表示时间的名词性词组“乾隆三十九年十二月”构成介宾结构，修饰谓语动词“乘”。这类介词“以”字是介入时间的，相当于古汉语中的介词“于”和现代的“在”。

但是，⑤的“以”只连接两个形容词，表示两种性质的联系，不能构成介宾结构，所以这类的“以”字是连词，相当于顺接连词“而”的用法，有“又”、“而”的意思。有时“以”字只连接两种行为动作，表示它们在时间上一一先后的联系，这类的“以”字同样是连词，相当于“而”，“有‘就、便’的意思，例如“齐因乘胜尽破其军，虏魏太子申以归。”（《史记·孙子吴起列传》）

⑥的“以”也是连接两种行为动作，但有表示目的的作用，可译为“来”、“以便”，称之为目的连词。

⑦的“以”处于主动词的位置，既非介词，也非连词，而是动词“认为”的意思。

介词“以”字可以用在谓语动词之前，它所构成的介宾结构充当动词的状语，已如上述。另外，它也可以用在谓语动词之后，它所构成的介宾结构作动词的补语，例如“何不试之以足？”（《郑人买履》），“为坛而盟，祭以尉首。”（《陈涉世家》）。这类用作补语的介宾结构，一般在翻译时要提到动词前面来，才符合现代汉语的习惯，象《郑人买履》中的这句话，可以译成“为什么不用你的脚试一试鞋子呢？”《陈涉世家》这句话，可以译成“筑起土台誓师，用将尉的头祭神。”

[为]
①苦 
为河伯
为河伯 娶妇，以故（因此）贫。（《西门豹治邺》）

②十余万人皆入睢水，睢水 
为之
为之 不流。（《史记·项羽本纪》）

③臣是凡人，偏在远郡，行将为人所并，岂足托乎？（《赤壁之战》）

④如今人方为刀俎，我为鱼肉，何辞为？（《鸿门宴》）

⑤霓为衣兮风为马，云之君兮纷纷而来下。（《梦游天姥吟留别》）

根据句意和介词特点，我们可以辨明①②③的“为”字以及④的最后一个“为”字，都是介词。④的最后一个“为”是语气词（①②的“为”读wěi；③的“为”和④的最后一个“为”，都读wéi）。至于④的前两个“为”字和⑤的“为”字都不符合介词特点，因为它们都作为句子的动词用，不能构成介宾结构去修饰什么谓语动词，它们都读wěi。

①的“为”与名词“河伯”构成介宾结构，修饰谓语动词“娶”。这类介词“为”字，经常和它后面的宾语一起放在谓语动词前面作状语，表示对象或目的，相当于现代的“替、给”。注意：这类“为”字的宾语有时可以省略，如“烦大巫妪为入报河伯”（《西门豹治邺》），我们根据句意和介词特点能够辨明“为”字后面省略了宾语“之”，译句时应把那省略了的宾语“之”所代的对象补出来。

②的“为”与代词“之”构成介宾结构，修饰谓语动词“流”。它所构成的介宾结构表示动作的原因，这类的“为”字相当于现代的“由于、因为”。

③的“为”与名词“人”构成介宾结构，修饰谓语动词“并”。这类“为”字的作用是介绍动词行为的主动者，相当于现代的“被”。它往往与“所”字呼应，组合成“为……所……”的被动句格式。

至于④中的前两个“为”字，正处于谓语动词的位置，起判断的作用，相当于现代的“是、象是”，第三个“为”字是介词，“何”是“辞”的前置宾语，句末语气词，常用在疑问句中，⑤中的两个“为”字，也正处于谓语动词的位置，相当于“作、当作”。由于这两个“为”字都不能构成介宾结构去修饰句子的谓语动词，所以能辨清它们都不是介词。

2.根据句意和“引进”、“介绍”的对象来确定介词词义或作用。

由于古代汉语中用作介词的字少，往往一个介词能表示现代汉语中好几个介词的词义。好在我们对现代汉语及口语有一定的正确的语言习惯和语感，可以根据句意和上下文把一个文言介词与现代汉语相当的词义表达出来。当然，我们如果能掌握文言介词译义的规律，就能讲得更准确、更有把握。例如：

于
于
①公与之乘，战于长勺。（《曹刿论战》）

②寡人之于国也，尽心焉耳矣。（《孟子·梁惠王上》）

③青，取之于蓝而青于蓝；冰，水为之而寒于水。（《劝学》）

④ 君幸于赵王，故燕王欲结于君。（《廉颇蔺相如列传》）

我们根据句意和介词特点来辨析，以上各句中的“于”字都是介词，因为它们都同名词（或代词）组成介词结构，修饰句中的谓语动词。所以说，“于”字是文言虚词中一个纯粹的介词。用在句子中的“乎”字，其作用大致相当于“于”字。

①的“于”与地名“长勺”结合，构成介宾结构，表示谓语动词“战”的处所。这类表示处所或时间的“于”，相当于现代的“在”或者“到”、“从”、“自”等。“于”字的这种用法，现代汉语仍然在沿用着。注意：这类表示处所的“于”字，在古代汉语的句子中有时省略，而此处所名词直接用在谓语动词或动宾词组之后作补语，如“令女居（于）其上，浮之（于）河中。”（《西门豹治邺》）我们在译句时，应补出这个省略了的“于”字。

②的“于”与名词“国”构成介宾结构，表示谓语动词“尽”的对象。这种用来介绍对象的“于”，相当于现代的“对、对于”或者“向”、“给”、“跟”等。这种用法，现代汉语也在沿用。注意：由于“于”字是个纯粹的介词，不要误把此句的“于”字当作动词看待。

③的第一个“于”字用法同①，相当于“从”。第二个和第三个“于”字引进比较的对象。这种引进比较的对象，表示比较的“于”字，一般只放在形容词谓语的后面，其意义相当于现代的“比”。

④的第一个“于”字，是引进动作行为的主动者，可译为“被”。它常同“受”、“见”呼应，构成“受……于……”“见……于……”的被动格式。第二个“于”字的用法同②，相当于现代的“跟、同”。

3.根据句意和古今词义的不同来辨析文言介词的词义或作用。

由于词义的发展变化，古代汉语的某些介词在现代汉语中已有了不同的讲法，我们在学习文言文时，要根据句意和古代汉语的通常词义细心辨析。如果用现代汉语的通常讲法来解释它，往往会出错。例如：

【因】

①廉颇闻之，肉袒负荆，因宾客至蔺相如门谢罪。（《廉颇蔺相如列传》）

②益州险塞，沃野千里，天府之土，高祖因之以成帝业。（《隆中对》）

③秦有余力而制其弊，追亡逐北，伏尸百万，流血漂橹；因利乘便，宰割天下，分裂山河。（《过秦论》）

④善战者因其势而利导之。（《史记·孙子吴起列传》）

⑤相如因持壁却立跪柱。（《廉颇蔺相如列传》）

⑥项羽即日因留沛公与饮。（《鸿门宴》）

⑦既而得其尸于井，因而化怒为悲。（《促织》）

⑧惠文、武、昭襄蒙故业，因遗策……（《过秦论》）

我们根据句意和介词特点，可以辨明①②③④中的“因”都是介词，因为它们能构成介宾结构，去修饰谓语动词。至于⑤⑥⑦中的“因”是连词，⑧中的“因”则是动词。值得注意的是，所有这些句子中的“因”，都不同于现代汉语作“因为”讲的“因”，“因”字在唐宋以前作“因为”讲的也有，但数量极少，需要注意辨析。
①的“因”，与名词“宾客”构成介宾结构，修饰谓语动词“至”，相当于现代的“通过”。
②的“因”，与代词“之”（代“益州”）构成介宾结构，修饰谓语动词“成”，相当于现代的“凭借”、“依靠”。
③的“因”，与用作名词的形容词“利”构成介宾结构，修饰谓语动词“宰割”、“分裂”，相当于现代的“趁着”。
④的“因”，与名词性词组“其势”构成介宾结构，修饰谓语动词“导”，相当于现代的“顺着”。
至于⑤⑥的“因”，都不具备介词的特点，只起连接的作用，是连词，相当于现代的“于是、就”。⑦的“因”，与“而”字连用，有“因之而”的意思，可译为“于是、从而”。⑧的“因”，处于谓语位置，是动词，当“沿袭”讲。

在古代汉语中，只有一个介词“与”，跟现代汉语的介词“与”的词义和用法相同，表示“同、跟”的意思；从“同、跟”的意思稍加引申，就是“替”、“给”的意思。例如：

①独卿与子敬与孤同耳。（《赤壁之战》）
②襄与吾祖居者，今其室十无一焉。（《捕蛇者说》）

③陈涉少时，尝与人佣耕。（《陈涉世家》）

句①中第二个“与”字和句②中的“与”字，都相当于现代的“同、跟”的意思；句③中的“与”字则跟现代的“给、替”相当。它们都是介词，而与句①中第一个“与”字不同。句①中第一个“与”字是连词，与现代的连词“和”相当。连词“与”连接平列的两个词语，连接的两项能换位置；介词“与”字的前后两个词语在句子结构上显然有主从之分，介词前后的两项不能换位置。

值得注意的是，古代汉语中的“与”经常用作动词，相当于现代的“给予”（如下面句①）、“结交”（如下面句②）、“参加”（如下面句③）、“对付”（如下面句④），它们的显著特点是处在句中的谓语位置上，而不是组成介词结构去作后面动词的状语。这一点需要认真仔细地鉴别。例如：

①吾与汝二万五千精兵，再拔一员上将，相助你去。（《失街亭》）

②齐人未尝赂秦，终继五国迁灭，何哉？与赢而不助五国也。（《六国论》）

③塞叔之子与师，哭而送之曰：“晋人御师必于毅。”（《左传·毅之战》）

④历阳侯范增曰：“汉易与耳！”（《史记·项羽本纪》）

此外，“与”字用在句末，是语气助词，表示疑问或感叹。如“可得闻与？”（《孟子·章·庄暴见孟子》）相当于现代的“吗”（有时相当于“呀”）。这个意义的“与”，后来写作“欤”。
\newpage